\section{The theory and its local cochain complex}\label{sec:setting}

Write $(\tau,z_1,z_2)$ for the coordinates of
$\mathbb R\times\mathbb C^2$ and
$\Omega=dz_1\wedge dz_2$. The superfield and its free differential are
\[
 A\in\Omega^\bullet(\mathbb R)\otimes
          \Omega^{0,\bullet}(\mathbb C^2)\otimes\g[1],
 \qquad
 DA=d\tau\,\partial_\tau A+\bar\partial A .
\]
The shift $[1]$ makes the total superfield odd. If $\kappa$ is the
invariant form, the action and the odd pairing are
\begin{equation}\label{eq:classical}
 \begin{split}
 S_{\mathrm{cl}}(A)
 &=\int\Omega\wedge
   \left(\frac12\kappa(A,DA)+\frac16\kappa(A,[A,A])\right),\\
 \omega(A,B)&=\int\Omega\wedge\kappa(A,B).
 \end{split}
\end{equation}
The bracket in the cubic combines the Lie bracket and exterior
multiplication. The pairing has degree $-1$.

\subsection{Locality and the sector being quantized}

Let $\Vht$ consist of constant-coefficient local functionals obtained
by applying finitely many holomorphic derivatives to each
superfield, taking the ordinary exterior product, and integrating
its degree-three component against $\Omega$. Functionals are taken
modulo constant-coefficient total derivatives. In particular no
real or antiholomorphic field derivatives, nor contractions of
external form indices, enter this space.

\begin{definition}\label{def:quantization}
A local formal quantization is an interaction
\[
 I+\sum_{k\ge1}\hbar^k I_k,\qquad
 I=\frac16\int\Omega\wedge\kappa(A,[A,A]),\qquad
 I_k\in\Vht,\quad \operatorname{gh}(I_k)=0,
\]
whose heat-kernel effective interactions solve the quantum master
equation at every positive scale, as formal power series.
Only positive field arity is at issue.
Equivariance means that $S/\hbar$ is invariant under
\[
 A\longmapsto\phi_g^*A,\qquad
 \hbar\longmapsto(\det g)^{-1}\hbar,\qquad g\in U(2),
\]
where $\phi_g$ acts on the two holomorphic coordinates.
\end{definition}

All superfield factors have degree one and the integral lowers
degree by three. Consequently a ghost-number-zero element of
$\Vht$ is cubic, and a ghost-number-one local obstruction is
quartic. Holomorphic derivatives do not change these degrees.
The equivariance condition on the order-$k$ cubic is
\[
 I_k(\phi_g^*A)=(\det g)^{k-1}I_k(A).
\]
For polynomial symbols this is equivalent to holomorphic
bidegree $(k,k)$ and invariance under $SL_2(\mathbb C)$.

\subsection{Momentum-conserving Taylor cochains}

The existing notation $C^0_{\g}$ denotes finite-jet
Chevalley--Eilenberg cochains of $\g[[z_1,z_2]]$, modulo the two
holomorphic total derivatives. An $r$-cochain is an alternating
multilinear function of $X_i e^{p_i\cdot z}$, polynomial in the
two-component holomorphic momenta, with
\[
 \sum_{i=0}^{r-1}p_i=0.
\]
This equation is the exact Fourier description of integration
by parts, not a restriction on a larger physical state space.
We use
\[
 [pq]=p_1q_2-p_2q_1,\qquad
 [Xe^{p\cdot z},Ye^{q\cdot z}]=[X,Y]e^{(p+q)\cdot z}.
\]
Here subscripts on a single momentum are coordinate indices.
The cochain differential, denoted by $d$, is
\begin{equation}\label{eq:CE}
 (d w)(x_0,\ldots,x_r)=
 \sum_{i<j}(-1)^{i+j}
 w([x_i,x_j],x_0,\ldots,\widehat{x_i},\ldots,
                      \widehat{x_j},\ldots,x_r).
\end{equation}
The hats omit those two arguments and $d^2=0$.
Full alternation means
$\alt w=\sum_{\sigma}\sgn(\sigma)\,w(x_{\sigma(0)},\ldots)$,
without division by a factorial.

An alternating cochain $w$ of arity $r$ gives the functional
$F_w=\frac1{r!}\int\Omega\wedge w(A,\ldots,A)$, with the
derivatives prescribed by its momentum symbol. Thus
\[
 C^{0,r}_{\g}\longrightarrow\Vht^{\,r-3}
\]
is an isomorphism of the indicated local complexes, with the
overall sign fixed below. This follows directly by polarization:
odd superfields enforce alternation, and total derivatives
enforce the sum of momenta. This dictionary concerns precisely
$\Vht$; it does not identify all translation-invariant local
functionals with this one summand.

The cochain bracket contracts one input of each cochain with
$\kappa^{-1}$ and sums the signed shuffles of the remaining
inputs. The contracted momenta are opposite, fixed by
conservation. If the two arities are $m$ and $n$, then
\begin{equation}\label{eq:cochain-bracket}
 \{a,b\}
 =\frac1{(m-1)!(n-1)!}
  \alt\left(
    a(x_1,\ldots,x_{m-1},X_e)
    \kappa^{ee'}
    b(X_{e'},x_m,\ldots,x_{m+n-2})
  \right).
\end{equation}
Here $X_e$ runs over any basis and $\kappa^{ee'}$ is the inverse
matrix of $\kappa$; the resulting contraction is basis independent.
For the classical cubic $s_0(X,Y,Z)=\kappa(X,[Y,Z])$,
\[
 \{s_0,w\}=-dw.
\]
The sign $\epsilon\in\{1,-1\}$ records the correspondence
$\{F_a,F_b\}=\epsilon F_{\{a,b\}}$. The sign $U\in\{1,-1\}$
records the common Wick sign multiplying the two-loop heat
integrals. We retain both in the existence argument.

For $J_1=F_j$ and $I_2=F_{i_2}$, the first two local equations are
\begin{equation}\label{eq:first-two-equations}
 dj=0,\qquad
 \epsilon\left(-d i_2+\frac12\{j,j\}\right)+o_2=0.
\end{equation}
Here $o_2$ is the two-loop cochain with its factor $\hbar^2$
removed. The one-loop local term vanishes, including one-loop
graphs containing holomorphically differentiated vertices.

\subsection{Heat normalization}

For the flat metric, use
\[
 D^*=-\iota_{\partial_\tau}\partial_\tau
      -4\sum_{i=1}^2\iota_{\partial_{\bar z_i}}\partial_{z_i}.
\]
The heat kernel normalized relative to $\Omega$ is
\begin{equation}\label{eq:heat}
 [K_t]_\Omega=
 \frac{\exp(-(\tau^2+|z_1|^2+|z_2|^2)/(4t))}
      {4(4\pi t)^{5/2}}\,
 d\tau\wedge d\bar z_1\wedge d\bar z_2,
 \qquad \int\Omega\wedge[K_t]_\Omega=1.
\end{equation}
The regularized propagator is
$P_\varepsilon^L=\int_\varepsilon^L D^*K_t\,dt$.
The common normalized scalar factor of a two-loop contact is
$-1/(64\pi^5)$.

Under the simultaneous dilation of coordinates by $r>0$ and
heat times by $r^2$,
\[
 r^*[K_{r^2t}]_\Omega=r^{-2}[K_t]_\Omega,\qquad
 r^*P_{r^2\varepsilon}^{r^2L}
 =r^{-2}P_\varepsilon^L.
\]
These are finite-cutoff identities. Their use for a limiting
local contact is justified by continuity of the decorated
graph distributions, not by replacing the regulated kernel
with a homogeneous expression prematurely.

